\section{Formulation}
\label{sec:method}

\begin{figure}[!t]
\centering
\input{\lfigdir/fig_overview.tex}
\caption{Overview. Planar surfels are the joint space of all sources. The depth network (black) and the
pixels (red) enter as sensors, each with a bias field on the surface, and each leaves its own information
record on the placement of every surfel through the same depth rows $\mathbf j_{uk}$. The records are made
honest, their disagreement is tested, and they are fused. Tags give the experimental stage that measures
each step (Section~\ref{sec:setup}): 0 the renderer and its Jacobian, A the depth sensor alone, B the
pixels, C both.}
\label{fig:overview}
\end{figure}

\subsection{The Surface as the Joint Space}
\label{sec:method:setting}
The input is a posed aerial video with frames $i=1,\dots,N$, images $\mathbf c_i$, poses
$(\mathrm R_i,\mathbf p_i)$ known up to a small error, and one depth map $\hat d_i$ per frame from a frozen
monocular network. The scene is a set of planar Gaussian surfels \cite{gs_2dgs}: surfel $k$ has center
$\boldsymbol\mu_k$, tangents $\mathbf t_{k1},\mathbf t_{k2}$, normal $\mathbf n_k$, scales, opacity and
color $\mathbf a_k$. Nothing is matched; both sources are expressed on the surfels, the joint space where they meet.

A surfel is a plane and so is the image, so the map between them is a homography. A point
$\boldsymbol\mu_k+\sigma_1\mathbf t_{k1}+\sigma_2\mathbf t_{k2}$ projects to
$\mathbf u\simeq\mathrm H_{ik}[\sigma_1,\sigma_2,1]^{\top}$ with
$\mathrm H_{ik}=\mathrm K\bigl[\mathrm R_i\mathbf t_{k1},\ \mathrm R_i\mathbf t_{k2},\
\mathrm R_i\boldsymbol\mu_k+\mathbf p_i\bigr]$, and two frames that see the surfel are related by
$\mathrm H_{jk}\mathrm H_{ik}^{-1}$, the homography induced by its plane. This is the only correspondence used, carried by the rasterizer \cite{gs_gsplat} through the
compositing weights $w_{uk}$ of pixel $u$; pose and surface enter the same matrix.

The \emph{placement} of surfel $k$ is $\boldsymbol\theta_k=(\nu_k,\tau_{k1},\tau_{k2})$: an offset along
the normal and two tilts, in meters and radians. For a ray of direction $\mathbf w_u$ (unit depth) that
meets the surfel at tangent coordinates $\boldsymbol\sigma_{uk}$, the depth changes with the placement by
the row
\begin{equation}
\mathbf j_{uk}=\frac{\partial d_{uk}}{\partial\boldsymbol\theta_k}
=\frac{1}{\mathbf n_k^{\top}\mathbf w_u}\bigl[\,1,\ \sigma_{uk,1},\ \sigma_{uk,2}\bigr].
\label{eq:depthjac}
\end{equation}

\subsection{Sources as Sensors with Bias Fields}
\label{sec:method:sensors}
Each source $s\in\{D,I\}$ is a sensor of the placement. Read alone, frame $i$ of source $s$ places surfel
$k$ at
\begin{equation}
\hat{\boldsymbol\theta}^{s}_{ik}=\boldsymbol\theta_k+\boldsymbol\beta^{s}_k+\boldsymbol\alpha^{s}_{ik},
\label{eq:sensor}
\end{equation}
where $\boldsymbol\alpha^{s}_{ik}$ differs between frames and $\boldsymbol\beta^{s}_k$ is a \emph{bias
field} on the surface, shared by every frame that sees surfel $k$. A frozen network repeats its mistake at the same
place, so $\boldsymbol\beta^{D}$ is large; non-Lambertian appearance and misfit give the pixels a bias
$\boldsymbol\beta^{I}$ of their own. Per frame the depth map also carries a scale $a_i$ and an offset $b_i$, and the pose an error
$\boldsymbol\xi_i\in\mathfrak{se}(3)$; together they form the frame nuisance $\mathbf z_i$.

The map is fit in two steps. First, depth-guided 2DGS is trained on a set of \emph{fitting} frames with
\begin{equation}
L=\sum_{i,u}\frac{\bigl(D_i(u)-a_i\hat d_i(u)-b_i\bigr)^2}{2\sigma_i^{2}}
+\frac{\lVert\hat{\mathbf c}_i(u)-\mathbf c_i(u)\rVert^2}{2\sigma_I^{2}},
\label{eq:loss}
\end{equation}
where $D_i=\sum_k w_{uk}d_{uk}$ and $\hat{\mathbf c}_i=\sum_k w_{uk}\mathbf a_k$ are rendered. The
geometry follows the depth term, the photometric term fits appearance and opacity, and poses and
$(a_i,b_i)$ are refined. The frozen surface is thus placed by the depth sensor, and its appearance never
saw the interleaved \emph{judging} frames rendered next.

\subsection{Per-Source Information Records}
\label{sec:method:records}
At the frozen map each source leaves a Gauss--Newton record on every surfel's placement.

\runin{Depth record.} Frame $i$ contributes
\begin{align}
\Lambda^{D}_{ik}&=\sum_{u}\frac{w_{uk}\,a_i^{2}}{\sigma_i^{2}}\,\mathbf j_{uk}^{\top}\mathbf j_{uk},&
\boldsymbol\eta^{D}_{ik}&=\sum_{u}\frac{w_{uk}\,a_i}{\sigma_i^{2}}\,\mathbf j_{uk}^{\top}r^{D}_i(u),
\label{eq:depthrecord}
\end{align}
with $r^{D}_i$ the depth residual of \eqref{eq:loss}. The weights enter linearly so that the record adds
over surfels.

\runin{Image record.} Moving surfel $k$ slides the point where a ray meets it: through
$\mathrm H_{ik}^{-1}$ the tangent coordinates move by $\mathbf m_{uk}\,\mathbf j_{uk}\boldsymbol\theta_k$,
with $\mathbf m_{uk}$ the tangential part of the ray in units of the surfel scales. The surfel's weight
changes with it and, through the transmittance, so do the weights behind it, so the pixel responds with
the contrast between the surfel's color and the color $\mathbf b_{uk}$ composited behind it:
\begin{equation}
\frac{\partial\hat{\mathbf c}_i(u)}{\partial\boldsymbol\theta_k}=\mathbf g_{uk}\,\mathbf j_{uk},\quad
\mathbf g_{uk}=-\,w_{uk}\,(\mathbf a_k-\mathbf b_{uk})\,(\boldsymbol\sigma_{uk}^{\top}\mathbf m_{uk}).
\label{eq:factor}
\end{equation}
The image Jacobian factors through the depth Jacobian: it is the chain rule through the homography, exact
to first order because the surfel is a plane, with surfel order and the color behind held fixed. Both
sources therefore speak about the same three coordinates. On a judging frame, the image record is
\begin{align}
\Lambda^{I}_{ik}&=\frac{1}{\sigma_I^{2}}\sum_{u}\mathbf j_{uk}^{\top}\mathbf g_{uk}^{\top}\mathbf g_{uk}\mathbf j_{uk},&
\boldsymbol\eta^{I}_{ik}&=\frac{1}{\sigma_I^{2}}\sum_{u}\mathbf j_{uk}^{\top}\mathbf g_{uk}^{\top}\mathbf r^{I}_i(u),
\label{eq:imagerecord}
\end{align}
with $\mathbf r^{I}_i$ the photometric residual. Surfels that share pixels are coupled by off-diagonal
blocks of the same form, kept sparse.

\runin{Nuisances.} The frame nuisance $\mathbf z_i$ is marginalized by a Schur complement,
$\Lambda^{s}_{ik}\leftarrow\Lambda^{s}_{ik}-\Lambda^{s}_{k z_i}\Lambda_{z_iz_i}^{-1}\Lambda^{s}_{z_i k}$,
where $\Lambda_{z_iz_i}$ pools every surfel in frame $i$ and the pose prior. A pose error moves every
surfel of a frame, a fault one place in every frame; a frame sees thousands of surfels, so a fault
barely moves $\mathbf z_i$. Records then add over frames, $\Lambda^{s}_{k}=\sum_i\Lambda^{s}_{ik}$, and leaving
out a block of frames is a subtraction.

\subsection{Honest Information}
\label{sec:method:honest}
Along a placement direction $\mathbf e$, let $\varrho_s$ be the share of the error variance of
source $s$ that its bias field carries,
$\varrho_s=\mathbf e^{\top}\Sigma^{s}_{\beta}\mathbf e/\mathbf e^{\top}(\Sigma^{s}_{\beta}+\Sigma^{s}_{\alpha})\mathbf e$.
\begin{proposition}\label{prop:law}
If the single-frame record of source $s$ is calibrated and its per-frame terms are independent, the fused
placement of $n$ equally informed frames has, along $\mathbf e$, an error variance
$\varrho_s+(1-\varrho_s)/n$ times that of one frame, whereas the summed record states $1/n$. The stated
variance is too small by $1+(n-1)\varrho_s$.
\end{proposition}
\begin{IEEEproof}
The fused error is $\boldsymbol\beta^{s}_k+\frac1n\sum_i\boldsymbol\alpha^{s}_{ik}$, of covariance
$\Sigma^{s}_\beta+\Sigma^{s}_\alpha/n$; the record states $(\Sigma^{s}_\beta+\Sigma^{s}_\alpha)/n$.
Project on $\mathbf e$ and divide.
\end{IEEEproof}
\noindent The honest information of a source saturates at its bias:
\begin{equation}
\tilde\Lambda^{s}_{k}=\bigl(\Sigma^{s}_\beta+(\Lambda^{s}_{k})^{-1}\bigr)^{-1}\ \preceq\ (\Sigma^{s}_\beta)^{-1},
\label{eq:saturation}
\end{equation}
which is covariance intersection \cite{fus_julier_ci} with the correlation estimated rather than bounded.
For the depth sensor, $\Sigma^{D}_\beta$ is the covariance that per-frame errors share across frames far
apart along the strip, measured on development data with exact depth and pooled by segment class. For the
pixels, a dispersion $\hat\varphi$ from frame replicates deflates $\Lambda^{I}$, and displacements injected
by moving surfels and re-rendering give the ratio of stated to realized information. Without truth, only
the difference $\boldsymbol\beta^{D}-\boldsymbol\beta^{I}$ is observable; on textured Lambertian ground,
where the injections show $\Sigma^{I}_\beta$ to be small, it is the depth bias.

\subsection{Testing the Sources Against Each Other}
\label{sec:method:test}
A fault is read where the depth sensor placed the surface, with appearance from the fitting frames, so
that appearance and opacity cannot be refit to hide it. On the judging frames the pixel pull is the score
$\mathbf s_k=\boldsymbol\eta^{I}_{k}$. If the depth-placed surface is off by $\boldsymbol\delta_k$, then
$\mathbf s_k\approx\Lambda^{I}_k\boldsymbol\delta_k$ plus noise of covariance $\hat\varphi\Lambda^{I}_k$.
A fault moves a surface, not a surfel, so the test is read on a \emph{segment} $q$: neighboring surfels
that a fault would move together. They are cut on the graph of surfels that share pixels, by agreement of
their planes and of frozen DINOv3 features \cite{feat_dinov3} lifted from the fitting frames only, so the
cut never sees the tested pixels. A segment is tested
against a basis $\mathrm B$ of fault shapes: each surfel alone ($\mathrm B=\mathrm I$), one rigid
motion per plane, or smooth deformations; a fault that crosses segments is caught at the next coarser
scale, with the scales corrected for multiple testing \cite{stat_holm}. With
$\mathbf s_q$ the stacked scores and
$\mathrm C_q=\hat\varphi\Lambda^{I}_q+\Lambda^{I}_q(\tilde\Lambda^{D}_q)^{-1}\Lambda^{I}_q$,
under the null that both sensors see one surface,
\begin{equation}
T_q=\mathbf s_q^{\top}\mathrm B\bigl(\mathrm B^{\top}\mathrm C_q\mathrm B\bigr)^{+}\mathrm B^{\top}\mathbf s_q\ \sim\ \chi^{2}_{r},
\label{eq:tension}
\end{equation}
with $r$ the rank of $\mathrm B^{\top}\mathrm C_q\mathrm B$, the number of fault directions in which the
pixels have stiffness. A depth error that every frame repeats cancels when depth maps are compared with
each other; it appears in $T_q$.

\subsection{Fusion and Scoring}
\label{sec:method:fusion}
The fused placement is one Gauss--Newton step from the frozen map,
$\Delta\boldsymbol\theta_k=\Lambda_k^{-1}\bigl(\tilde{\boldsymbol\eta}^{D}_{k}+\hat\varphi^{-1}\boldsymbol\eta^{I}_{k}\bigr)$
with $\Lambda_k=\tilde\Lambda^{D}_{k}+\hat\varphi^{-1}\Lambda^{I}_{k}$ its stated information, where
$\tilde{\boldsymbol\eta}^{D}_{k}$ is the depth pull scaled as in \eqref{eq:saturation}. On segments where
$T_q$ exceeds the $\chi^2_r$ quantile of a stated false-alarm rate the depth record is left out, and in
directions without pixel stiffness the surface is reported as unsupported. A map is scored by the
\emph{surface log-score}, the mean over surfels of $\tfrac12 z_k^2-\tfrac12\log\det\Lambda_k$, where
$z_k^2=\boldsymbol\epsilon_k^{\top}\Lambda_k\boldsymbol\epsilon_k$ and $\boldsymbol\epsilon_k$ is the
placement error against exact points, tilts in meters at the surfel edge. It is a proper scoring rule
\cite{stat_gneiting_scoring}: it cannot be improved by inflating or deflating the stated information.
